% ============================================================
% ESQUELETO DE MANUSCRITO — Bioengineering (MDPI)
% Creado: 03-ago-2026 — DECISION FINAL de revista (reemplaza a POI/SAGE)
%
% COMO USAR ESTE ARCHIVO (importante, leer antes de pegar en Overleaf):
%
% 1. Este archivo NO es autocontenido -- no trae la clase mdpi.cls ni la
%    carpeta Definitions/ (logos, estilos, macros de cada revista MDPI).
%    Esos archivos NO se pudieron descargar aqui. El flujo correcto es:
%      a) Abrir la plantilla oficial verificada en Overleaf:
%         https://www.overleaf.com/latex/templates/mdpi-article-template/fcpwsspfzsph
%         (click "Open as Template" -- crea un proyecto nuevo con
%         template.tex + carpeta Definitions/ ya completa)
%      b) Reemplazar el contenido de "template.tex" de ese proyecto con
%         el contenido de este archivo (de \documentclass en adelante)
%      c) MDPI actualiza su plantilla varias veces al ano y la version
%         de Overleaf casi siempre esta un poco atras de mdpi.com/authors/latex
%         -- si algo no compila, revisar ahi antes de asumir que es un
%         error de este archivo.
%
% 2. Overleaf tiene envio directo a revistas MDPI: boton "Submit" en el
%    proyecto -> "Submit to an MDPI journal" -- empaqueta y redirige al
%    sistema de envio de MDPI. Util para cuando el manuscrito este listo.
%
% 3. Los comandos de front-matter (\Title, \Author, \abstract, \keyword,
%    \authorcontributions, etc.) SI estan verificados contra la
%    plantilla oficial de Overleaf -- no son una suposicion.
%
% 4. APC confirmado: CHF 2700 (Bioengineering, 2025) -- cubierto por la
%    universidad segun lo indicado en la sesion del 03-ago-2026, no es
%    limitante para esta eleccion.
%
% 5. Contenido: igual que en el intento anterior (POI/SAGE, ver
%    plantilla_overleaf_POI.tex -- queda como referencia, no se borra,
%    pero la decision de revista final es esta), tomado de
%    docs/manuscrito/metodos_introduccion_borrador.md, ya sin
%    Kinovea/IMU de Alessandro/cita al paper de conferencia.
%
% 6. Estilo de referencias de MDPI: no se confirmo el comando exacto de
%    \bibliography en este ciclo -- revisar el template.tex real
%    descargado de Overleaf, que ya trae la configuracion correcta.
% ============================================================

\documentclass[bioengineering,article,submit,moreauthors]{Definitions/mdpi}

% ------------------------------------------------------------
% TITULO
% ------------------------------------------------------------
\Title{Multi-Subject Functional Validation of a 3-DOF Gait Simulator for
Transtibial Prosthesis Testing Using Inertial Motion Capture}

% [PENDIENTE: nombres, afiliaciones y orden de autoria -- decidir con CRediT statement]
\Author{Autor Uno $^{1}$, Autor Dos $^{1}$, Autor Tres $^{2}$, Autor Cuatro $^{3,}$*}

\AuthorNames{Autor Uno, Autor Dos, Autor Tres, Autor Cuatro}

\address{%
$^{1}$ Departamento de Ingenier\'ia, Pontificia Universidad Cat\'olica del Per\'u (PUCP), Lima, Per\'u\\
$^{2}$ [PENDIENTE: afiliaci\'on]\\
$^{3}$ [PENDIENTE: afiliaci\'on]}

\corres{Correspondence: [PENDIENTE: correo]; Tel.: [PENDIENTE]}

% ------------------------------------------------------------
% ABSTRACT -- MDPI pide estructura Background/Methods/Results/Conclusions, maximo 200 palabras
% ------------------------------------------------------------
\abstract{Background: [PENDIENTE]. Methods: [PENDIENTE]. Results: [PENDIENTE --
depende de datos de sujetos nuevos y del piloto del IMU en la plataforma,
ver CLAUDE.md]. Conclusions: [PENDIENTE].}

\keyword{gait simulator; transtibial prosthesis; inertial motion capture; multi-subject validation; prosthetics testing}
% MDPI: separar keywords con punto y coma, no coma -- confirmar limite maximo

\begin{document}

% ============================================================
\section{Introduction}
% ============================================================
% [PENDIENTE: cuerpo completo -- necesidad clinica con literatura
% externa, sin autocita al paper de conferencia. Ver nota en
% metodos_introduccion_borrador.md sobre por que no se cita.]

Transtibial prosthesis testing increasingly relies on gait simulators to
provide controlled, repeatable loading conditions for functional and
mechanical evaluation without requiring continuous human-subject
involvement. A central open question for any such simulator is whether
trajectories programmed from a single reference subject generalize to
the gait patterns of individuals who did not contribute to that
programming --- a question that is rarely addressed with quantified,
multi-subject evidence in the literature on prosthesis gait simulators.
The present study addresses this gap directly: we assess whether a
3-degree-of-freedom (horizontal, vertical, sagittal) motor-driven gait
simulator reproduces gait kinematics captured from multiple subjects
independent of its calibration trajectory, using an inertial motion
capture system whose validity for lower-limb kinematics --- including
specifically in transtibial prosthesis users --- is already established
in the literature \cite{PENDIENTE_iSen_transtibial}, and we quantify ---
rather than narrate --- the sources of vertical force overestimation
inherent to a motor-driven platform.

% ============================================================
\section{Materials and Methods}
% ============================================================

\subsection{System Description}
% Condensado, autocontenido -- no remite al paper de conferencia (no citado en este articulo)

The gait simulator reproduces a pre-recorded gait trajectory across
three independent degrees of freedom --- horizontal displacement,
vertical displacement, and sagittal rotation --- each driven by its own
motor and transmission (chain-and-sprocket or direct drive, depending on
the axis). Trajectories are stored as CSV files, uploaded via a
Raspberry Pi, and executed open-loop by an ESP32 microcontroller; the
present study does not modify this operating principle --- no
closed-loop control or algorithmic trajectory generation was introduced
in this cycle (see Discussion, Future Work). % [PENDIENTE: calibracion del actuador vertical, bloqueada por integracion Raspberry Pi-ESP32]

\subsection{Instrumentation}
% Un solo instrumento en todo el protocolo (revisado 03-ago-2026)

Kinematic and kinetic data were collected using two instruments:

\begin{itemize}
\item \textit{Inertial motion capture (STT-IWS/iSen).} Used throughout
the protocol to (i) capture new subjects' natural gait, (ii) re-capture
the original reference subject to regenerate the simulator's default
trajectory, and (iii) measure the simulator's own platform output
(sensor mounted directly on the platform). Segment/platform inclination
angles are derived from the sensor's onboard orientation estimate
relative to gravity, positive above horizontal and negative below, the
same sign convention used throughout this project. No in-house
cross-instrument validation was performed in this cycle; instrument
validity for lower-limb gait kinematics --- including specifically in
transtibial prosthesis users --- is supported by previously published
validation studies rather than re-derived \cite{PENDIENTE_iSen_transtibial}.
% [PENDIENTE: confirmar que el sensor exporta orientacion cruda respecto
% a la gravedad y no solo el angulo articular calibrado del protocolo de
% marcha por defecto -- resultado del piloto en curso. Ver tambien
% docs/literatura/postprocesado_datos_crudos_IMU.md sobre el corte de
% filtro que hay que recalibrar para la velocidad real del simulador.]
\item \textit{Force platform (AMTI).} Used to record the vertical
ground reaction force (Fz) exerted by the simulator against a fixed
load cell, sampled at 1000 Hz.
% [PENDIENTE: confirmar si se reportan las 6 columnas completas (Fx,Fy,Fz,Mx,My,Mz) o solo Fz]
\end{itemize}

\subsubsection{Vertical Force Overestimation: A Three-Stage, Decoupled Explanation}

Because the moving assembly is not a single rigid body --- a stationary
electronic control board is mounted on one side of the frame, while
three independently actuated axes (horizontal, vertical, sagittal) each
carry their own kinematic chain on the other --- vertical force
overestimation is decomposed into three independent, sequentially
corrected sources rather than a single inertial term:

\begin{enumerate}
\item A fixed \textit{vertical datum offset}, calibrated once against
an independent static load reference (not against any subject
trajectory reported in Results) and held constant across all trials.
\item \textit{Command-to-encoder tracking fidelity} per axis, isolating
mechanical/control error (chain or sprocket backlash, insufficient
motor holding torque under reaction load) from the remaining sources.
\item An \textit{axis-wise inertial correction}, using the mass that
actually accelerates with each motor according to the real mounting
order of the kinematic chain --- not the total assembly mass ---
benchmarked against published prosthetic gait GRF data.
% [PENDIENTE: fuente de m_eje sin resolver -- CAD vs. renunciar al termino y declararlo limitacion]
\end{enumerate}

\subsection{Experimental Protocol}
% [PENDIENTE: depende de si llego la aprobacion de etica y del numero final de sujetos]

\subsection{Statistical Analysis}
% [PENDIENTE: prosa -- formulas listas en
% docs/planificacion/plan_trabajo_5_semanas_articulo_Q2.md. Incluir:
% RMSEnorm, Pearson r (con nota de cautela en curvas monotonas), %
% puntos dentro de +-1 SD, ICC(3,1) [Koo & Li 2016] especificando el
% modelo exacto, SPM1D no parametrico por permutacion [Nichols & Holmes
% 2002; Pataky et al. 2015]. NO mencionar Bland-Altman (sin uso este
% ciclo). Encuadrar exactitud citando ISO 5725 -- ver
% docs/literatura/normas_ISO_relevantes.md.]

% ============================================================
\section{Results}
% ============================================================
% Orden -- ver docs/planificacion/propuesta_articulo_Q2.md secc. 5

\subsection{Tracking Fidelity per Subject}
% [PENDIENTE: depende de datos de sujetos nuevos]

\subsection{Representativeness of the Default Trajectory}
% [PENDIENTE: depende de datos de sujetos nuevos]

\subsection{Repeatability and Between-Subject Variability}
% [PENDIENTE: depende de datos de sujetos nuevos]

\subsection{Vertical Force: Raw, Corrected, and Literature Benchmark}
% [PENDIENTE: depende de la calibracion del offset vertical, bloqueada por integracion Raspberry Pi-ESP32]

% ============================================================
\section{Discussion}
% ============================================================
% [PENDIENTE: interpretacion, limitaciones explicitas (tamano de
% muestra), subseccion "Sources of error, decoupled", trabajo futuro
% (lazo cerrado, generacion algoritmica de trayectorias, IMU de
% Alessandro y celda de carga para el segundo articulo)]

% ============================================================
\section{Conclusions}
% ============================================================
% [PENDIENTE]

% ------------------------------------------------------------
% DECLARACIONES -- comandos verificados contra la plantilla oficial de Overleaf
% ------------------------------------------------------------
\authorcontributions{% [PENDIENTE: taxonomia CRediT -- decidir quien hizo que]
For research articles with several authors, a short paragraph specifying
their individual contributions must be provided using the CREDIT taxonomy.}

\funding{% [PENDIENTE]
This research received no external funding. / [PENDIENTE si aplica financiamiento real]}

\institutionalreview{% [PENDIENTE: depende de aprobacion del comite de etica]
Ethical review and approval [PENDIENTE] by the [PENDIENTE: nombre del comite], protocol code [PENDIENTE], date of approval [PENDIENTE].}

\informedconsent{Informed consent was obtained from all subjects involved in the study.} % [PENDIENTE: confirmar]

\dataavailability{% [PENDIENTE: decidir alcance -- repo completo publico vs. solo scripts de analisis con datos de ejemplo, ver propuesta_articulo_Q2.md secc. 7]
}

\conflictsofinterest{The authors declare no conflict of interest.} % [PENDIENTE: confirmar]

\acknowledgments{[PENDIENTE]}

\abbreviations{Abbreviations}{% [PENDIENTE: completar segun se usen en el texto]
The following abbreviations are used in this manuscript:\\
\noindent
\begin{tabular}{@{}ll}
GRF & Ground Reaction Force\\
IMU & Inertial Measurement Unit\\
RMSE & Root Mean Square Error\\
ICC & Intraclass Correlation Coefficient\\
SPM & Statistical Parametric Mapping\\
\end{tabular}}

% Referencias ya identificadas para incluir (ver docs/literatura/):
%   - validacion_instrumentos_IMU.md: validez iSen/STT-IWS (confirmar cita exacta antes de usar)
%   - validacion_instrumentos_IMU.md: IMU en usuarios de protesis transtibial (Sensors 2023, DOI 10.3390/s23031738)
%   - literatura_GRF_protesica.md: benchmark de vGRF en marcha protesica
%   - normas_ISO_relevantes.md: ISO 5725; Wu et al. 2002 (convencion ISB)
%   - postprocesado_datos_crudos_IMU.md: filtrado/calibracion de IMU
%   - Koo & Li (2016) -- ICC(3,1)
%   - Nichols & Holmes (2002); Pataky et al. (2015) -- SPM1D no parametrico
%   - Kadaba et al. (1989) -- CMC
% [PENDIENTE: el comando exacto de bibliografia de MDPI no se confirmo
% en este ciclo -- usar el que ya trae el template.tex real descargado
% de Overleaf, normalmente algo como \bibliography{references} o
% \reftitle{References} seguido de \bibliography]

\end{document}
